\documentclass{osucourse}
\course{5591H}
%% Imported from the sheet Math 5591H shares with 5590H; the sections below are common to both.

\begin{document}

\CatalogDescription
\PrerequisitesAndExclusions

\begin{Text}
Vary, for example:  Abstract Algebra, 3rd edition, by Dummit \& Foote, published by Wiley, ISBN: 9780471433349  Algebra, by M. Artin  Topics in Algebra, by I. Herstein
\end{Text}

\begin{TopicsOutline}
5590H:

\topicsub{1. Integers, unique factorization; congruences, Euler function.}

\topicsub{2. Groups, subgroups, homomorphisms and isomorphisms, normal subgroups, quotient groups, permutation groups, cyclic groups, Cauchy Theorems, Sylow's Theorems; direct products, fundamental theorem for finite Abelian group; G-sets.}

\topicsub{3. Rings, subrings, ideals, morphisms, polynomial rings, prime and maximal ideals.}

\topicsub{4. Commutative rings, factorization theory, Euclidean rings, principal ideal rings, unique factorization domains, Gauss' lemma; illustrations in the integers of quadratic number fields.}

\topicsub{5. Modules over commutative rings, submodules, quotients and direct sums; fundamental theorem for modules over principal ideal domains.}

5591H:

\topicsub{1. Vector spaces (as a special case of modules); linear maps and matrices, canonical forms, dual spaces.}

\topicsub{2. The theory of determinants.}

\topicsub{3. Bilinear and quadratic forms; inner product and unitary spaces; principal axis theorem.}

\topicsub{4. Fields, algebraic and transcendental (extensions), existence of closure (over countable fields), tests for polynomial irreducibility; normality, separability, field automorphisms.}

\topicsub{5. Galois theory, the subgroup-subfield correspondence theorem, group theory interrelations; extensions of finite fields, cyclotomic extensions.}

\topicsub{6. Solvable groups and solvability by radicals.}
\end{TopicsOutline}

\end{document}
